\chapter{Deviations from the protocols}
\label{chap:09-ecarts}

A protocol written in advance protects a result only if every later change is visible. This chapter
describes the rule we followed, what a reader can check of it and what rests on our word. It then
gives the complete list of deviations from the three protocols and discusses those that weigh most
on the results.

\section{The principle}
\label{sec:ecarts-principe}

The three protocols date from 8 November 2026: the survey protocol (\file{study/PROTOCOL.md}),
written before any paper was read, the simulation protocol (\file{sim/SIM\_PROTOCOL.md}) and the
protocol of the GRPO strand (\file{sim/llm/LLM\_PROTOCOL.md}), both written before any training.
Each one fixes the question, the population or the runs, the measures and the analyses. Every later
change is appended at the bottom of the file, in a numbered list of \enquote{deviations}. The body
of the protocol is never rewritten.

Each deviation records three things: its date, its reason and what we had already seen when we
decided it. The third matters most. A rule fixed before a result is seen cannot have been chosen to
suit that result. A rule fixed afterwards may have been. The deviations therefore say
\enquote{BEFORE} or \enquote{AFTER}, in capitals, with the exact object of the comparison:
before any training, after the PPO and SAC results but before any CrossQ result, after the values
at step 100 were displayed but before any gain was computed. Three rules complete the arrangement.
No run is removed. A run that crashes is relaunched once with the same random seed, then the
incident is recorded. Once an analysis has been announced, its result is published whatever it is.

The three protocols list 66 deviations: 29 for the survey, 20 for the simulations, 17 for the GRPO
strand. \Cref{tab:ecarts} summarises all of them. \Cref{fig:frise-ecarts} places them in time,
from 8 to 21 November 2026.

\section{What rests on our word}
\label{sec:parole}

A reader can recompute every number from the committed results (\cref{chap:08-logiciel}) and check
every survey decision against the quotation and the page it rests on. A reader cannot check, unaided,
the order in which things were written.

The repository became public on 15 November 2026. Its first commit, dated 9 November at 06:07,
already holds results next to the protocols. For everything before 15 November, the order of writing
rests on our word and on the commit dates, which their author could have forged. A critical review of
the public repository also found several \enquote{before}/\enquote{after} statements that were wrong
in their details. The corrections are dated and published: deviation 14 of the simulations,
deviation 15 of the GRPO strand. They show that the rule was applied with errors. The reader learns
this before deciding how far to trust the other statements.

The repository offers three supports for that trust. The evaluation logs of every run are committed,
including those of stopped runs when they were not overwritten. The attempt ledger lists every
training launch. The provenance of the survey readers (instructions, model, dates, papers written)
is reconstructed from the transcripts and published in \file{study/readers/} (survey deviation 28).

\section{The attempt ledgers}
\label{sec:registres}

The simulation protocol states that a run that crashes is relaunched once and the incident recorded.
Until 15 November, no ledger kept track of every launch. The reconstructed ledger,
\file{sim/logs/attempts\_reconstructed.csv}, rests on the creation dates of the logs, the outputs of
the launchers, the committed versions and our session logs. It gives one line per attempt:
algorithm, task, random seed, attempt number, start, end, outcome, score seen before the relaunch and
evidence. Simulation deviation 14 (e) draws up the balance: 292 launches for the 250 training runs of
the first phase, 42 attempts stopped (only 2 of them declared) and none crashed. Every stop hit all
the runs in progress at once. No run was relaunched or discarded after its final score was known. For
PPO PyTorch and PPO JAX, the relaunches reproduce the evaluation points of the stopped attempts
identically. For SAC and CrossQ, the stopped curves were overwritten and do not allow this check.

Since 15 November, the launcher \file{queue.sh} keeps one log per attempt, never overwritten. It
appends each attempt to the ledger \file{sim/logs/attempts.jsonl}: start, end, exit code, SHA-256
fingerprint of the result. The watchdog writes into the ledger every job it stops. A stop by signal
is no longer relaunched automatically. The stopped GRPO runs are recorded there too.

\section{The list of deviations}
\label{sec:liste-ecarts}

\gentable{ecarts}

\fig{frise-ecarts}{The protocols and their dated deviations, from 8 to 21 November 2026. One line per
protocol (survey, simulations, GRPO strand). Each mark is a deviation, placed at its date. Its shape
shows whether it was written before or after the results it concerns.}

\section{Deviations of the survey}
\label{sec:ecarts-sondage}

\paragraph{Harmonisation (deviations 1 to 6, 15 and 19)} The decision rules of the protocol left
borderline cases open. Deviations 1 to 6, recorded during reading and before any percentage, fix
the most frequent ones: X4 requires a described setting for the baselines, X7 an amount of compute,
Q5 a link that leads to code. Each change is recoded from the quotation, in a harmonisation file
(\file{overrides.json}, \file{overrides\_v2.json}) that gives the reason for every modified value.
Agreement between readers is always computed on the raw readings, before harmonisation. Deviation 15
corrects a mistake of ours, seen after the percentages: the harmonisation rules had been applied
only to the 2026 papers. The same rules, with no new rule added, were then applied to the twelve
2018 papers of the calibration. Deviation 19 applies the harmonisation of grid v2 to both years,
after an audit of the raw readings and before any percentage: 80 changes, each with its reason. For
X3, X4 and X7, no new rule was added, so that both years are read with the same grid. Borderline
cases stay coded as the reader read them.

\paragraph{Grid v2 and its planning (deviations 11, 14, 16 and 17)} Grid v1 measured only part of the
flaws. Grid v2, Y1 to Y12, was written on 9 November 2026 before any v2 reading (deviation 11), its
decision rules the same day (deviation 14). They were made more precise after a pilot on three papers
(deviations 16 and 17), still before the main coding. The work plan that precedes these deviations,
\file{study/planification\_v2.md}, was kept as it was and made public on 15 November 2026. For each
flaw, it states which v2 criteria are meant to move the verdict out of the \enquote{to be measured}
state.

\paragraph{The enlargement (deviation 12)} The power audit of the first draw shows that twelve 2018
papers and fourteen classical RL papers from 2026 are too few to detect a gap of thirty points.
Deviation 12, written before any draw and before any of the added papers was read, raises each of
the two samples to 50 papers and fixes the main comparison: classical RL in 2026 against our own
reading of 2018, with the same grids. The papers of the 2026 queue go through a title screen that
sets aside titles presumed to concern large language models. \Cref{chap:10-limites} discusses its
effect.

\paragraph{The verdict rule (deviation 20)} Of all the survey deviations, this one weighs most on the
conclusions. It was decided after the first results of the enlarged reading had been seen. Until
then, a criterion counted as changed when the gap held against Pineau's figures. The secondary
comparison shows that our grid reads 2018 differently from her: 55\% against 90\% for the
hyperparameters of the method (Q3), 33\% against 10\% for the held-out evaluation (Q7).
Comparing 2026 with her figures would mix the change in practices with the change in reading. Since
this deviation, a verdict rests only on the main comparison of deviation 12, fixed before any of the
added papers was read. The comparisons with Pineau are still shown, as context. The change of rule is
declared as decided after the results. The comparison it retains had been fixed before them.

\paragraph{The human review dropped (deviations 13 and 23)} Deviation 13 planned for Warith Harchaoui
to reread ten 2026 papers alone, blind, on seven criteria. Deviation 23 dropped that plan on
10 November 2026, before any rereading. As a result, no human reread the papers. Every decision can
still be checked against its quotation. The folder \file{study/relecture\_humaine/} is kept as a
trace of the plan.

\paragraph{Analyses written in response to review (deviations 21, 22, 24 to 29)} After the results,
several analyses were added to answer objections. Each was written before its own computation and
published whatever it showed: the robustness cuts and the test of the cost thesis (deviation 21), the
checklists declared by the authors (deviation 22), the title filter (deviation 24), the yearly
trajectory and its second reading (deviations 25 and 26), the mechanical search against the
\enquote{no} decisions (deviation 27), the provenance of the readers (deviation 28) and the removal of
the thirteen papers partly read by another model (deviation 29). Deviation 21 removed one reading
from the document: the thesis \enquote{Pineau was heard where it cost little} fails the correlation
test that the deviation announced (\cref{chap:04-sondage-resultats}).

\section{Deviations of the simulations}
\label{sec:ecarts-sim}

\paragraph{The budget (deviation 1)} SAC and CrossQ train for 200{,}000 steps instead of 300{,}000, a
decision taken before the first training to keep the compute within reach. PPO keeps its million
steps. This deviation creates the comparison at unequal budgets, which deviation 12 (iii) then
corrects with a reading at equal budget.

\paragraph{B3, written along the way (deviation 8)} Experiment B3, which asks whether the measure
(final score, best score, mean of the curve) changes the ranking, was written on 10 November 2026,
after the PPO and SAC results. It trains nothing: it rereads the curves already recorded. Deviation 8
dated it \enquote{BEFORE any CrossQ result}. Deviation 14 (b) corrects that account, since seven
CrossQ runs had finished when B3 was written. Our session logs, which are not public, show that no
CrossQ score had been read. The exact wording becomes \enquote{before looking at any CrossQ result}.
Deviation 7, which fixes the details of E2, E4 and E5, receives the same correction.

\paragraph{The descriptive ranking under modified physics (deviation 11)} The most quoted reversal of
the blog article, PPO JAX first on \task{Hopper} in the training physics and last in the modified
physics, comes from a complement decided after the E7 results had been seen, without a test
(\file{analysis/e7\_rank.py}). The document presents it as such. The graded re-evaluation of
deviation 12 (ii), written before any new computation, then put this reversal to the test: PPO JAX is
no longer first on \task{Hopper} in 10 of the 12 modified physics.

\paragraph{The consolidation (deviation 12)} Written on 15 November 2026, before any new training and
any new evaluation, this deviation adds three analyses with the same code and the same settings:
twenty random seeds per pair (160 training runs), the re-evaluation of every policy under masses and
friction modified by degrees, plus the reading of PPO at 200{,}000 steps, the budget of SAC and
CrossQ. The results at ten random seeds remain the main results. \Cref{chap:06-experiences-resultats} gives both
next to each other and states every divergence.

\paragraph{The random seed of the bootstrap (deviation 13)} While the public repository was being
prepared, six scripts did not reproduce their intervals identically. rliable ignores its
\texttt{random\_state} argument and draws from NumPy's global generator. Each script now seeds this
generator at the start. Two runs now give files identical byte for byte. The point estimates and the
verdicts do not change. The bounds move at the level of Monte Carlo noise, by at most 16 points on an
interquartile mean of 2{,}423 for E6.

\paragraph{Seven corrections after a critical review (deviation 14)} A hostile review of the public
repository established seven points, none of which changes a verdict.
(a) The commit that described the dose script as \enquote{written before the dose results} was
wrong: 30 of the 160 re-evaluations already existed. The script applies a rule written before any
computation, with no free parameter.
(b) The counts of finished CrossQ runs in deviations 7 and 8 were wrong.
(c) The protocol announced Stable-Baselines3 2.6.0. Every training ran with version 2.9.2, installed
before the first training.
(d) PPO and TRPO collect in batches of 2{,}048 steps, so each block of 50{,}000 steps is in fact
51{,}200 steps: the PPO and TRPO models have 1{,}024{,}000 steps, against 200{,}000 for SAC and
CrossQ.
(e) The undeclared relaunches, described in \cref{sec:registres}.
(f) Six scripts read every file present. With random seeds 10 to 19, E6 would have gone from ten to
twenty random seeds without saying so. Each script now names its random seeds.
(g) The ledger \file{attempts.jsonl} and the per-attempt logs, from 15 November on.

\paragraph{Organising the compute (deviations 9, 10, 15 to 17 and 19)} These deviations record stops
by the watchdog, relaunches, the move to eight cores and then to two, plus a suspension tied to the
Mac's battery, lifted the same morning. None changes a result: each run uses a single thread and
depends only on its own random seed.

\paragraph{Computing the additions (deviation 20)} Written on 21 November 2026, after every training
had finished and before any computation on random seeds 10 to 19, this deviation fixes a defect of
\file{analysis/e5.py}, which took the mean of the first ten runs as the \enquote{correct} ordering
even with twenty random seeds. It then lists the computations, all part of \texttt{make analysis}. Two
of them, the descriptive ranking and the Bonferroni correction at twenty random seeds, were not in
deviation 12: they apply rules already written and are published whatever they show.

\section{Deviations of the GRPO strand}
\label{sec:ecarts-llm}

\paragraph{The plan reduced, then restored (deviations 6 and 14)} Deviation 6, decided before any
retained run, reduces the plan to 100 steps and three random seeds, sets the learning rate to
$3 \times 10^{-6}$ without a full pilot and splits the update into sub-batches. Deviation 14, written
on 15 November 2026 before any new run, adds random seeds 3 and 4. Of all the decisions, this
restoration changed the reading of the strand most: Qwen's gain with the true reward, clear at three
random seeds, is no longer clear at five (\cref{chap:07-modeles-de-langue}).

\paragraph{The crashes and their cause (deviations 7 to 11)} The first retained runs crashed on
non-finite probabilities or ran out of GPU memory. Deviations 7 and 8 set a relaunch rule, decided
before any result of these runs was seen. Deviation 9 moves the runs to the night. Deviation 10 finds
the cause: a GPU computation interrupted by macOS returns with corrupted values. It adds the
safeguards of \file{grpo.py}. It also resets the failure count to zero, since the past crashes came
from the machine. This change of rule was decided after the crashes were seen, at a time when no
result of these runs existed. Deviation 11 moves the remaining Qwen runs to the CPU. It is the
origin of the confounding between hardware and random seed discussed in \cref{sec:llm-materiel}.

\paragraph{Corrections after review (deviation 15)} Written on 15 November 2026, after the results,
this deviation corrects the breakdown of the nine failures of 9 November and the date of one run,
pins the exact revision of both models in \file{grpo.py}, documents the stops of 15 November and
commits to comparing the CPU runs alone, since the device changes between Qwen's random seeds.

\paragraph{The run that fails twice (deviation 16)} This deviation was written on 17 November 2026,
after the step-100 values of the seven other new runs had been displayed during a progress check,
but before any gain or interval at five random seeds was computed. The Llama run with the randomly
drawn reward and random seed 4 fails twice on the gradient safeguard. Under deviation 14, this second
failure is a result. The configuration is summarised over its four finished runs, with a sensitivity
analysis for the missing gain. The deviation also plans one CPU run outside the protocol, whose
result is kept apart.

\paragraph{The order of the CPU run (deviation 17)} With the two-core ceiling of simulation
deviation 19, the Llama CPU run starts after all the other computations. It has no effect on the
results.

\section{What these deviations say about the work}
\label{sec:ecarts-bilan}

Most deviations have one of three causes: the compute available on one laptop, our own
mistakes, objections raised in review. The deviations decided after the results are declared as
such. The three that weigh most on a conclusion are the verdict rule of the survey (deviation 20),
the descriptive ranking under modified physics (simulation deviation 11) and the reset of the GRPO
failure count (deviation 10). For the first, the retained comparison had been fixed before the
results. The second was put to the test by an analysis written in advance. The third affects no
published number, since no crashed run had a result. A reader who rejects these three decisions can
redo every computation without them from the committed files.
